\subsection{One finite five-point equation supplies the retained genus-zero action}

\begin{result}{A.4.37}{Proved}{An explicit characteristic five-point level produces invertible finite braid and sphere actions}
Let $m\geq2$ be a specified integer and retain all arities
$n\leq10$.  For the ordinary pure braid group $P_n$ put
\[
 K_n=\kappa_m(P_n),\qquad
 C=(m!)^{(m!)^{45}},\qquad
 Z_5=\Gamma_{0,5}/\kappa_C(\Gamma_{0,5}).                \tag{FPA.1}
\]
All $K_n$ are literal intersections of kernels of
homomorphisms to $\mathfrak S_m$.
For an actual gentle element $F=(\lambda,f)\in\widehat{\mathrm{GT}}_0$,
consider its ordinary pentagon word
\[
\begin{split}
 p(f)={}&f(A_{12},A_{23}A_{24})f(A_{13}A_{23},A_{34})\\
 &\quad\cdot
 \bigl(f(A_{23},A_{34})f(A_{12}A_{13},A_{24}A_{34})
                                      f(A_{12},A_{23})\bigr)^{-1}
                  \quad\in\widehat P_4.
\end{split}                                                    \tag{FPA.2}
\]
Here $A_{ij}=\sigma_{j-1}\cdots\sigma_{i+1}\sigma_i^2
\sigma_{i+1}^{-1}\cdots\sigma_{j-1}^{-1}$, and the map
$P_4\to\Gamma_{0,5}$ kills its central full twist.
Every element of the locus defined by the finite equation
\[
 \{F\in\widehat{\mathrm{GT}}_0: p(f)=1\text{ in }Z_5\}
                                                               \tag{FPA.3}
\]
induces well-defined automorphisms of the finite
unitary parenthesized-braid operad whose arity-$n$ pure
loop group is $P_n/K_n$.  The induced actions on
\[
 S_n(m)=\Gamma_{0,[n]}/\kappa_m(\Gamma_{0,n}),
                    \qquad 3\leq n\leq10,             \tag{FPA.4}
\]
fix the label permutations and commute with the actual
labelled forgetting maps, on their specified stabilizers.
They are the standard finite braid formulas with the
same pair $(\lambda,f)$.  Every prescribed finite quotient
of $\Gamma_{0,[n]}$ of order at most $m$ is retained by
(FPA.4), by an actual quotient map.  No full pentagon
identity or full $\widehat{\mathrm{GT}}$ lift is assumed
for an element of (FPA.3).
\end{result}

\begin{proof}
\par\medskip\noindent\textbf{The finite operad and every insertion map.}\quad
Objects in arity $n$ are the parenthesized permutations
of the $n$ labels.  A morphism is the usual braid with
the appropriate permutation.  At every object quotient
its pure loop group by $K_n$; transporting $K_n$ between
objects is independent of the chosen braid, because
$K_n$ is characteristic in $P_n$ and normal in $B_n$.
This gives a finite groupoid $\mathcal B_m(n)$.

For any homomorphism $u:G\to H$ one has
\[
                     u(\kappa_m(G))\subseteq\kappa_m(H):
                                                               \tag{FPA.5}
\]
compose $u$ with each homomorphism $H\to\mathfrak S_m$.
Every pure insertion, deletion or relabelling of a braid
is such a homomorphism on the appropriate pure loop
group.  Inserting a loop in either input of an operadic
composition therefore carries its $K$ subgroup into the
output $K$ subgroup.  A change of the comparison path
only conjugates that subgroup, which is already normal.
Thus every actual braid insertion and every deletion
descends to these finite groupoids.  They form a
finite unitary operad truncated at arity ten; its
composition maps are the specified ordinary cabling maps.
The pure group $P_n$ has $\binom n2\leq45$ named pair
generators, so the diagonal construction gives
\[
               |P_n/K_n|\leq(m!)^{(m!)^{\binom n2}}\leq C.
                                                               \tag{FPA.6}
\]
No assertion of functoriality of an unspecified finite
quotient is used here.

\par\medskip\noindent\textbf{The exact five-point cap of every inserted test.}\quad
Take any one of the ordinary four-block insertions
$P_4\to P_n$ and its finite composite $q:P_4\to D$,
where $D$ is its actual image in $P_n/K_n$.
Let $z$ be the central full twist of $P_4$ and
$n_z=\operatorname{ord}(q(z))$.
Use the literal abelian character
\[
 \chi:P_4\longrightarrow\mathbb Z,\qquad
 \chi(A_{12})=1,\quad\chi(A_{ij})=0\ ((i,j)\ne(1,2)).
                                                               \tag{FPA.7}
\]
The abelianization of $z$ has coefficient one on every
pair generator, so $\chi(z)=1$.
Form the actual finite join and cap
\[
 \widetilde D=\operatorname{im}(q,\chi\bmod n_z),\qquad
 \zeta=(q(z),1),\qquad D^\flat=\widetilde D/\langle\zeta\rangle.
                                                               \tag{FPA.8}
\]
The element $\zeta$ is central of order $n_z$, so this
is an actual quotient of $P_4/\langle z\rangle=\Gamma_{0,5}$
and
\[
                    |D^\flat|\leq |D|\leq C.           \tag{FPA.9}
\]
Its regular permutation representation, enlarged by fixed
letters to degree $C$, is one of the maps defining
$\kappa_C(\Gamma_{0,5})$.
Consequently (FPA.3) makes the image of $p(f)$ trivial
in $D^\flat$.
Each factor of (FPA.2) is a value of a commutator proword,
so $p(f)$ has zero abelianization in $\widehat P_4$ and
in particular zero $\chi$ coordinate.
If $(q(p(f)),0)=\zeta^j$, the second coordinate says
$j=0\pmod{n_z}$, and hence $q(p(f))=1$.
This proves the needed pentagon in $D$ itself, not merely
in its central quotient.

For clarity, the five homomorphisms evaluating the
coefficient in the basic four-strand pentagon are
\[
\begin{array}{c|cc}
 &x&y\\ \hline
 1&A_{12}&A_{23}A_{24}\\
 2&A_{13}A_{23}&A_{34}\\
 3&A_{23}&A_{34}\\
 4&A_{12}A_{13}&A_{24}A_{34}\\
 5&A_{12}&A_{23}.
\end{array}                                                    \tag{FPA.10}
\]
They are obtained by the three strand doublings and
the two exterior insertions.  Composing any row with
any four-block cabling into at most ten strands is the
actual finite evaluation just capped in (FPA.8).
Thus the single characteristic quotient in (FPA.1)
retains every required inserted pentagon, including
all relabelled and nested instances.

\par\medskip\noindent\textbf{The presentation and the new finite endofunctor.}\quad
Use the discrete parenthesized-braid presentation by the
binary object, the braiding $\tau$ and the associator
$\alpha$, with pentagon, the two hexagons, and unit
relations; see Fresse, \emph{Homotopy of Operads and
Grothendieck--Teichm\"uller Groups}, Part I, Theorem~6.2.4,
pp.~212--218 \cite{Fresse2017}.  That theorem is a presentation of ordinary
groupoids, before any Malcev or profinite completion.
Its proof constructs a morphism in arity $n$ using only
insertions in arities at most $n$, so it applies to this
finite truncation as well.
Send the objects to themselves and set
\[
       \tau\longmapsto\tau^\lambda,\qquad
       \alpha\longmapsto\alpha f(A_{12},A_{23}).        \tag{FPA.11}
\]
Composition is read from right to left.  The image of
$A_{12}$ at the left-associated three-point object is
$A_{12}^\lambda$, and that of $A_{23}$ is
$f^{-1}A_{23}^\lambda f$.
The two paths of the basic associator pentagon then
have precisely the two products displayed in (FPA.2):
the double insertion at the first, middle and last
positions gives the rows $2,4,1$ of (FPA.10), respectively.
The previous paragraph proves equality in every required
finite target.
The two hexagon images are the ordinary gentle
three-cycle and antisymmetry equations in the three-strand
group; their insertion maps are actual homomorphisms,
so those equations hold here for every actual gentle
pair.  The unit equations are $f(x,1)=f(1,y)=1$, which
hold because $f$ is a commutator proword.
These are the exact finite images of the named equations,
not additional shadow conditions.  The discrete
presentation therefore supplies a functor from the
ordinary parenthesized-braid operad to $\mathcal B_m$.
Profinite values of $f$ are evaluated in these finite
loop groups, so this is also the continuous functor
from their full completions with the same generator
images.

We now verify that this new functor factors through
$\mathcal B_m$ itself.  This does not follow just from
functoriality of the original cablings.
On an arity-$n$ pure loop group it is a homomorphism
\[
                         P_n\longrightarrow P_n/K_n.
                                                               \tag{FPA.12}
\]
Compose (FPA.12) with each of the defining diagonal
coordinate maps $P_n/K_n\to\mathfrak S_m$.
Every resulting homomorphism $P_n\to\mathfrak S_m$
kills $K_n$ by its definition.  The coordinates are
jointly faithful, so (FPA.12) kills $K_n$.
Thus the functor descends to an endofunctor of the same
finite groupoids.  Since the ordinary functor already
commutes with the insertion and unit maps, its descent
is an endomorphism of the finite truncated unitary
operad.

\par\medskip\noindent\textbf{Invertibility from a finite genuine gentle record.}\quad
Put
\[
 R=10!\,C,\qquad E=\widehat F_2/\kappa_R(\widehat F_2),
                       L=2|E|.                        \tag{FPA.13}
\]
Every group in the retained braid groupoids has order
at most $R$ after including its permutation part.
Any evaluation of a two-generator proword in such a
group factors through $E$, by the regular-representation
argument.  The exponent of every such two-generated
image therefore divides $|E|$; $\lambda\bmod L$ also
retains all half-exponents $(\lambda-1)/2$ used by the
hexagons.
The actual gentle group has a finite group image given
by the record
\[
                  (\lambda\bmod L,\ f\bmod\kappa_R(F_2)).
                                                               \tag{FPA.14}
\]
Indeed $\kappa_R(F_2)$ is characteristic, and the
normalized gentle automorphism
$\beta_F(x)=x^\lambda$, $\beta_F(y)=f^{-1}y^\lambda f$
acts on $E$.  With this composition convention the
pair law is
\[
 (\lambda,f)(\mu,g)
             =(\lambda\mu,\ f\,\beta_F(g)).            \tag{FPA.15}
\]
This formula descends to (FPA.14), and the actual image
is a finite group.  It is not replaced by the set of
all formal tuples in the displayed finite sets.

Composition of the endofunctor (FPA.11) with another
such endofunctor has exactly (FPA.15) as its generator
record.  On $\tau$ the powers multiply; on $\alpha$
its existing coefficient is transformed by
$x\mapsto x^\lambda$, $y\mapsto f^{-1}y^\lambda f$,
which gives the second entry of (FPA.15).
This calculation takes place in the already constructed
finite category and does not assume a pentagon for an
inverse pair.
The particular record of $F$ has finite order $d$ in
(FPA.14).  Its $d$-th endofunctor power consequently
fixes $\tau$ and $\alpha$, and therefore every morphism
by the presentation.  Thus that power is the identity,
and the endofunctor is invertible.
For an explicit uniform bound, the record group has
order at most $\varphi(L)|E|$, so
$(\varphi(L)|E|)!$ is an allowable common power.
This proves invertibility without a full pentagon lift.

\par\medskip\noindent\textbf{The sphere quotient and its preserved universal kernel.}\quad
The finite braid action sends the full twist $z_n$ to
$z_n^\lambda$.  To verify this directly, decompose a
parenthesized object into two blocks: its full twist is
the product of the two internal full twists and the
double interchange of the blocks.  Those three factors
commute.  Operadic compatibility sends the double
interchange to its $\lambda$-th power; induction on the
blocks gives the assertion.  The result is independent
of reassociation because the full twist is central.
For the standard nested prefix twists one has
\[
 y_n=z_nz_{n-1}^{-1}
       =\sigma_{n-1}\cdots\sigma_1^2\cdots\sigma_{n-1}.
                                                               \tag{FPA.16}
\]
The two factors commute and both are powered by
$\lambda$, so $y_n$ is powered by $\lambda$.
The usual spherical braid relator is conjugate to
$y_n$, and the other mapping-class relator is $z_n=1$.
Their normal closure is therefore preserved.  This
gives the finite spherical mapping-class action on
the quotient of $B_n/K_n$ by those two relators.
Its invertibility follows from the same finite power
argument, or from the already invertible braid action
and preservation of a finite normal subgroup.

Finally project to (FPA.4).  On the pure sphere group
the resulting map is a homomorphism
$\Gamma_{0,n}\to\Gamma_{0,n}/\kappa_m(\Gamma_{0,n})$.
Composing with each defining coordinate to
$\mathfrak S_m$ proves that it kills the source
$\kappa_m(\Gamma_{0,n})$, exactly as in (FPA.12).
Thus it induces an automorphism of (FPA.4).
Its characteristic pure kernel is normal in the full
mixed group; label permutations are unchanged because
$\lambda$ is odd and associator coefficients are pure.
Deleting labelled strands is the unitary operad map,
which commutes with (FPA.11).  The same coordinate
argument applies to the induced sphere-group maps.
The literal quotient chain is
\[
\begin{gathered}
 B_n/K_n\longrightarrow
 B_n/(K_n\langle\!\langle z_n,y_n\rangle\!\rangle)
 \longrightarrow S_n(m),\\
 B_n\longrightarrow\Gamma_{0,[n]}.
\end{gathered}
                                                               \tag{FPA.17}
\]
The composite of the top chain with $B_n\to B_n/K_n$
is the composite of the bottom capping map with the
specified quotient of the sphere group.  Each arrow
intertwines the just-constructed finite endomorphisms;
this is the quotient square, including the pure-sphere
kernel which can be larger than the image of $K_n$.
This proves the claimed finite forgetting compatibility.
A prescribed quotient of order at most $m$ kills the
pure $\kappa_m$ by its actual regular representation,
so its source map factors through (FPA.4), as asserted.
The executed program
\path{work/tower_code/finite_pentagon_action_words.py}
independently verifies all three strand-doubling rows
in (FPA.10), every binary-block full-twist decomposition
and its commutations, and the literal conjugacy of the
two sphere words for $2\leq n\leq10$, using the faithful
free-group Artin action.  It also checks the full-twist
abelian coordinates in (FPA.7).  The source, output and
JSON braid words are retained.
\end{proof}

\paragraph{Nested levels for the finite subgroup refinements.}
Start with a required finite ten-point output of order
at most $m_0$ and use its literal map from $S_{10}(m_0)$.
Apply A.4.34 or A.4.35 to the actual restriction of that
universal output to the named colored subgroup, and
induce the resulting finite image to the full mixed
group by its finite permutation cosets.  Let $m_1$ be
any of the proved upper bounds for the resulting finite
image, increased to be at least $m_0$.
Then $S_{10}(m_1)$ maps onto that actual intermediate
refinement and onto $S_{10}(m_0)$.
Using (FPA.1) with $m=m_1$ supplies the finite actions
at both levels, and the generator formulas prove that
these actions commute with their literal projection.
The intermediate kernel need not be declared invariant.
This is the precise way the finite fusion and centralizer
statements enter the common five-point input.
The remaining application must still retain its full
source shadow image and perform the demanded census;
no enumeration of all genuine shadows is claimed here.
